\section{Introduction}\label{sec:introduction84}
A measurement with a classical output may retain a quantum advantage
when its input is correlated with a reference.  The question considered
here is quantitative: how well can two ordered measurements be
separated with at most $N$ uses, and which features of their effects
determine the dependence on $N$?  The answer must accommodate both
strictly positive effects and singular measurements.  Positivity of
each effect alone is not the relevant constraint: the effects sum to
the identity, and their boundary geometry is coupled by this equation.

\subsection{The model and the principal results}
An ordered $k$-outcome measurement on $\C^d$ is a tuple
$E_j\succeq0$ with $\sum_jE_j=I_d$.  One call is the channel
\[
 \mathcal M_E(\rho)=\sum_{j=1}^k\tr(E_j\rho)|j\rangle\langle j|.
\]
The input is consumed and the output is only the classical label;
no residual device system is available.  A tester may prepare inputs
entangled with its own reference, retain quantum memory, and choose
subsequent controls using previous labels.  It uses the same control
rule under both hypotheses.  Its public stopping time is bounded
by $N$.  We define $D_N(E,F)$ to be the supremum of the unhalved
trace distance between its two final states.  Thus $0\le D_N\le2$.
Nonadaptive testers may have entangled inputs and retained references;
``nonadaptive'' does not mean ``product input.''


\subsection{Complete profiles and hard quadratic budgets}
The principal refinement retains the full acquisition allocation,
rather than only its largest permissible group.  For
$\boldsymbol n=(n_1,\ldots,n_\ell)$ put
$N=\sum_r n_r$ and $V=\sum_r n_r^2$.
Theorem~\ref{thm:allocation89} proves the three local orders
\begin{equation}\label{eq:frontprofile89}
 \min\{1,\sqrt N s\},\qquad
 \min\{1,\sqrt V s\},\qquad
 \min\{1,Ns\}
\end{equation}
for regular tangents, coherent tangents and support openings,
respectively.  They hold both for independent complete groups and
for prescribed-profile reset-constrained feedback.  In the latter
class each entire fresh probe, reference and within-block memory
is independent of all old receiver memory conditional on the
classical history.  Preparations use one common rule defined even
at histories null under one hypothesis.  The two conditional old
receiver states need not agree.

Theorem~\ref{thm:quadraticbudget89} allows the reserved group sizes
to depend on history.  If $\sum n_r\le N$ and $\sum n_r^2\le Q$
on every formal terminal path, where $N\le Q\le N^2$, then
optimization over this reset budget class replaces $\sqrt V s$
by $\sqrt Q s$ in the coherent row of \eqref{eq:frontprofile89}.
The other rows are unchanged.  These are hard budgets, not expected
costs.  The finer finite bound for a specified policy is
\[
 \min\{2,\,2s(\alpha\sqrt{Q_\pi}+\sqrt{rN_\pi})\},
\]
where the second term retains the full quadratic tube remainder.
A specified policy need not attain a lower of this order: it can
forget its outputs or spend its large budget on uninformative paths.
The matching assertion concerns the best tester in the budget class.

The lower combines unequal GHZ groups using one bounded classical
readout determined by public signal sizes, not by the unknown
remainder.  Each $m_r\kappa s^2$ channel error is paid before
amplification.  The upper uses a conditional fidelity induction on
subnormalized receiver operators; it never asserts unconditional
product outputs after feedback.  Neither proof incurs a logarithmic
loss from grouping allocations by size.  The sum-of-squares principle,
GHZ enhancement and fidelity methods have substantial antecedents
\cite{Hyllus78,Toth78,HMNW88,YF76}.  The additional statement is the
finite measurement-tube comparison, uniform in every complete
profile and hard quadratic budget, with remainder-independent lower
controls.

Reset width is neither memory dimension nor the entanglement depth
of probe marginals.  Proposition~\ref{prop:memorycomparison89}
proves that the two-call measure-and-reprepare testers of
\cite[Definition~23]{OZNPQ89} form a proper subset of our unit-reset
class at the level of tester operators.  The witness uses retained
receiver references and a final joint readout, not coherent import
into a later probe.  It does not assert an optimal-score separation
for every channel ensemble.  Proposition~\ref{prop:approxreset89}
separately quantifies the effect of a certified uniform channel
error from reset; the certificate is an explicit hypothesis,
not a consequence of a policy graph.

\subsection{Finite neighborhoods and entanglement width}
Fix $E$, a nonzero Hermitian tuple $H$ with $\sum_jH_j=0$, and
$\Lambda\ge0$.  For a Hermitian tuple write
$\|R\|_\Sigma=\sum_j\|R_j\|_{\mathrm{op}}$.  Our neighborhood
consists of all legal measurements $F$ such that
\begin{equation}\label{eq:fronttube87}
 \|F-E-sH\|_\Sigma\le\Lambda s^2.
\end{equation}
No curve or smooth Kraus representation is specified, and the
remainder may vary arbitrarily with $s$.  Put
\[
 P_j=\operatorname{supp}E_j,\quad Q_j=I-P_j,\quad
 J_j=Q_jH_jQ_j,\quad
 \Gamma_E(H)=\frac i2\sum_j[P_j,H_j],\quad
 \mathcal W_E=\sum_jP_j\mathcal H_dP_j.
\]
Here $\mathcal H_d$ is the real Hermitian matrix space.  The
one-sided cone is $J_j\succeq0$; its lineality space is $J_j=0$.
The positivity condition is standard semidefinite tangent geometry
\cite{BCS87}.  The coupled normalization, support-range test and
finite-use consequences are proved below.  In particular,
\[
 H\in\operatorname{Ran}\mathcal C_E
 \quad\Longleftrightarrow\quad
 J_j=0\ (1\le j\le k),\qquad \Gamma_E(H)\in\mathcal W_E.
\]

Let $D_N^{[b]}$ allow independent groups of complete probe--reference
systems, each using at most $b$ calls, with arbitrary entanglement
within a group.  All inputs precede the calls; all outputs may be
processed jointly afterwards.  Definition~\ref{def:blocks87} includes
public mixtures and records.  The endpoints are independent
probe--reference acquisition at $b=1$ and arbitrary entangled
parallel acquisition at $b=N$.  Theorem~\ref{thm:widthlaw87} proves
\begin{equation}\label{eq:frontwidth87}
 D_N^{[b]}(E,F)\asymp_{E,H,\Lambda}
 \begin{cases}
 \min\{1,\sqrt N s\},&H\in\operatorname{Ran}\mathcal C_E,\\
 \min\{1,\sqrt{Nb}\,s\},&J_j=0\ \text{for all }j,
                        \quad\Gamma_E(H)\notin\mathcal W_E,\\
 \min\{1,Ns\},&\text{some }J_j\ne0.
 \end{cases}
\end{equation}
The comparison is uniform over all integers $1\le b\le N$,
all sufficiently small $s>0$, and every $F$ in
\eqref{eq:fronttube87}.  The constants and the interval depend on
the fixed $E,H,\Lambda$, not on $b,N$ or the unknown remainder.
The same lower preparations and readouts work throughout the tube.
They may depend on $E,H,s,N,b$ and are not common-learning advice.

These rows separate three mechanisms.  A regular tangent has only
square-root accumulation even with full adaptive control.  A
coherent tangent gains from entanglement across device inputs.
A support opening instead creates a classically impossible outcome,
so independent repeated inputs already give linear accumulation.
The new parallel construction encodes GHZ blocks before all calls
and recovers each logical system only at the end.  Its remainder is
$O(ms^2)$ on a block of $m$ calls; repeated independent blocks and
a finite Bernoulli inequality then give the middle row.  The upper
uses Bures comparison within each entangled block and fidelity
tensorization only between independent blocks.

In particular Corollary~\ref{cor:parallel87} yields
\[
 D_N^{\mathrm{par}}(E,F)\asymp_{E,H,\Lambda}D_N(E,F).
\]
This is equality of local orders, not equality of exact distances,
perfect-discrimination thresholds, or fixed-pair error exponents.
The earlier product converse is not applied to correlated inputs.
GHZ metrology, entanglement-depth bounds and the Kraus-span correction
mechanism have established antecedents
\cite{Hyllus78,Toth78,ZJ84,KL84}.  The present result is a
finite effect-neighborhood theorem, with explicit remainder control
and uniformity also in the permitted block width.

The neighborhood is not presumed nonempty for arbitrary $\Lambda$.
Choose $a>0$ with $\|H_j\|\le a$ and a common support floor
$E_j\succeq\lambda P_j$, $\lambda>0$.  For
\[
 c=1+2a^2/\lambda,\qquad
 s_* =\min\{1,\lambda/(2a),1/a\},\qquad
 \Lambda_*=ck(1+2k),
\]
Proposition~\ref{prop:nonempty87} gives a legal member
$(E_j+sH_j+cs^2I)/(1+kcs^2)$ for every $s\in(0,s_*]$
whenever $\Lambda\ge\Lambda_*$.  This is a sufficient, not
minimal, allowance.  For smaller $\Lambda$ all rate statements
quantify only over the legal members.  The displayed curve is a
rational function of $s$; its coefficients are Gaussian-rational
when the effects, directions and auxiliary bounds are so represented.
The exact certificate concerns legality and remainder control,
not the small interval of the discrimination theorem.

The finite-neighborhood result includes one-sided $C^2$ curves with
nonzero right derivative and controlled higher-order families
$E+t^qH+O(t^{2q})$.  In the latter case
Corollary~\ref{cor:widthpower87} gives $\sqrt{Nb}\,t^q$
in the coherent branch.  A mere leading coefficient is not sufficient:
the retained scalar example has mixed order
$\min\{1,\sqrt N t^a+Nt^b\}$, $1\le a<b<2a$.
No general zero-jet classification or arbitrary-pair midpoint
metric equivalence is inferred.

\subsection{Classical feedback and receiver memory}
Theorem~\ref{thm:resetwidth88} proves the same three rates for a
larger class.  A fresh block may now contain arbitrary adaptive
quantum control.  Between blocks, the receiver may process all its
quantum memory jointly and send classical messages determining the
next complete probe--reference preparation.  Conditional on each
message history, that new preparation must be independent of the
receiver memory; no coherent subsystem of an earlier block enters
the new acquisition.  This reset condition, rather than the
entanglement depth of probe marginals, defines the resource.

For a history-dependent allocation the finite upper depends on the
worst-path quadratic cost $Q=\sup_h\sum_r n(h_r)^2\le bN$.
A conditional fidelity lemma treats arbitrary incoming receiver
states at each history.  It replaces the invalid multiplication of
unconditional output fidelities after feedback.  Fresh adaptive
blocks obey the finite estimate $s(\alpha n+\sqrt{rn})$ in Bures
distance, retaining the $nrs^2$ channel remainder.  The resulting
trace upper is $2(\alpha+\sqrt r)s\sqrt Q$, capped at two.
The corrected GHZ lower from the independent-block theorem already
belongs to this larger class.  Its $m\kappa s^2$ remainder and its
readout independent of the unknown tube remainder are unchanged.

In particular, unit reset width admits classical feedback and
unrestricted receiver quantum memory but retains the square-root
local order on coherent tangents.  A single unrestricted block
admits all adaptive protocols.  This statement concerns finite local
orders under the stated reset condition, not exact equality of
strategy distances or fixed-pair error exponents.  Channel-Bures
bounds and classical-feedback discrimination are established
antecedents \cite{HMNW88,SHW87}; the present theorem supplies the
fixed-tube, hard-budget interpolation.

\subsection{The covariance and support structure}
Our complete-body upper bound uses a positive, concave covariance
$\mathcal C_E$ on zero-sum Hermitian tuples.  Its factorization through
a horizontal projection gives a finite-use bound valid through changes
of rank and zero effects.  With $M=(E+F)/2$ and $H=F-E$,
\begin{equation}\label{eq:frontupper84}
 D_N(E,F)\le\min\left\{2,2\sqrt{k+1}
  \sqrt{N\langle H,(\mathcal C_M+N^{-1}\operatorname{Id})^{-1}H\rangle}
                              \right\}.
\end{equation}
This is a global upper certificate, not a claimed formula for the
optimal adaptive distance.  The normalization includes the exact
binary Sylvester restriction of Theorem~\ref{thm:closedcovariance83}.

The support description determines its entire kernel.  If
$P_j=\operatorname{supp}E_j$ and
$\mathcal W_E=\sum_jP_j\mathcal H_dP_j$, then
Theorem~\ref{thm:supportkernel84} parametrizes all null directions
and proves
\[
 \dim\ker(\mathcal C_E|_{\mathcal T})
 =d^2-\dim\mathcal W_E+\sum_j(d-\operatorname{rank}E_j)^2.
\]
This distinction between support span and the mere effect span is
essential for effects of rank greater than one.  In particular,
a nonprojective measurement can have a nontrivial covariance kernel;
the earlier projective characterization concerns a zero operator,
not one null direction.

This kernel has a finite-use interpretation on every fixed unitary
orbit $E_j(t)=e^{itG}E_je^{-itG}$.  Apart from the stationary case,
Theorem~\ref{thm:orbittrichotomy84} proves, uniformly in every integer
$N\ge1$ and all sufficiently small $|t|$,
\begin{equation}\label{eq:frontorbit84}
 D_N(E,E(t))\asymp_{E,G}
 \begin{cases}
 \min\{1,\sqrt N|t|\},&G\in\mathcal W_E,\\
 \min\{1,N|t|\},&G\notin\mathcal W_E.
 \end{cases}
\end{equation}
The second lower bound uses a corrected logical qubit and the explicit
finite-angle estimate \eqref{eq:finiteanglelower84}.
The support span is precisely the Hermitian Kraus-product span of
this measurement channel.  The metrological criterion is credited
to Zhou--Jiang \cite{ZJ84}; the present formulation identifies it
with the covariance kernel and proves the finite-pair trace estimates
directly.  This is an orbit theorem on the complete body, not a
full-boundary volume or minimax theorem.

We also replace a coordinate-fixed ridge by a public scalar-measurement
covariance.  Transporting its probability vector through the same
output channel gives a regularized data-processing inequality and
exact invariance under reversible label splitting
(Theorem~\ref{thm:transportedridge84}).  The original fixed Euclidean
ridge is not asserted to have this property.

\subsection{Nonunitary curves and certified controls}
The support description also gives a local classification without a
unitary-orbit assumption.  For a fixed two-sided $C^2$ measurement curve
with $H=E'(0)\ne0$, set
\[
 \Gamma_E(H)=\frac i2\sum_j[\operatorname{supp}E_j,H_j].
\]
Theorem~\ref{thm:curveclassification85} proves
\[
 D_N(E(0),E(t))\asymp_{E(\cdot)}
 \begin{cases}
 \min\{1,\sqrt N|t|\},&\Gamma_E(H)\in\mathcal W_E,\\
 \min\{1,N|t|\},&\Gamma_E(H)\notin\mathcal W_E.
 \end{cases}
\]
This holds for every finite $N$ on a two-sided interval depending on
the fixed curve.  The supports may change at second order.  An analytic
horizontal surrogate supplies a $\sqrt N|t|+O(Nt^2)$ upper in the
range branch; a support-complement code supplies a logical channel
with an explicit $O(t^2)$ remainder in the other branch.  Both are
finite inequalities rather than inferences from asymptotic Fisher
information.  The metrological exponents and correction mechanism
remain attributed to \cite{ZJ84,KL84}; the new formulation works
directly with $C^2$ effects even when an exact Kraus family is not
assumed differentiable at a rank opening.

Theorem~\ref{thm:controlstability85} bounds the loss from uniformly
certified errors in preparation, encoding, recovery and readout.
The allowed per-cycle error scales with $|t|\Delta$, where $\Delta$
is the logical gap.  This is a conditional stability guarantee for
known-pair controls, not an executed implementation or a gate-synthesis
complexity theorem.  The constants and interval may degenerate with
support eigenvalues, tangent size, curvature and the support residual.
Zero first derivative does not imply stationarity for a general curve.
The global arbitrary-pair covariance bound remains one-sided.

\subsection{Learning and descriptions on balanced interiors}
The operational consequences are complemented by exact description
and common learning laws on the balanced interior
$I_d/(2k)\preceq E_j\preceq3I_d/(2k)$.  Its real affine dimension
is $p=(k-1)d^2$.  Theorems~\ref{thm:multioutcomeentropy82} and
\ref{thm:multioutcomecode82} give
\[
 \log_2\mathcal C_{N,k}(\delta)
 =\frac p2\log_2N+p\log_2(1/\delta)+O(p\log(k+1)),
 \qquad 0<\delta<1/128.
\]
The common learner has the joint order
\[
 M^*_{d,k}(N,\delta,\eta)
 =\Theta_k\!\left(N\delta^{-2}[d^2+d\log(1/\eta)]\right)
\]
for fixed $k$, $d,N\ge1$, $0<\delta\le2^{-24}$ and
$0<\eta\le1/8$.  The displayed constructive upper has a $k^3$
factor; the lower does not match growing $k$.
Theorem~\ref{thm:affinerepair83} supplies polynomial, search-free
statistical legalization, while the entropy-optimal public dictionary
remains a separate finite, potentially exhaustive construction.
The lower theorem permits arbitrary coherent adaptive training,
whereas the upper uses fresh acquisitions and collective processing
of completed outputs.  Pair-dependent discrimination witnesses do
not serve as a common learning design.

\subsection{Resource conventions and reading order}
Unknown-device calls, future-use horizon, transmitted index length,
public dictionary reconstruction, classical arithmetic, trusted-control
synthesis, and physical execution are distinct resources.  A query
bound here holds on every record, with a defined legal fallback.
Public dictionaries and their parameters are independent of the
target; target-dependent private advice would count as transmitted
information.  Exact rational certificate evaluation and balanced
legalization have polynomial bit complexity.  Neither statement
asserts polynomial synthesis of an arbitrary collective readout or
of the whole entropy-optimal dictionary.  The orbit recovery is an
ideal-control comparison construction, not an executed common learner.

Sections~\ref{sec:covariance83}--\ref{sec:ridge84} develop the global
upper, support kernel, orbit classification and transported ridge.
Sections~\ref{sec:finiteoutcome82}--\ref{sec:affinerepair83} give the
balanced-interior learning and coding consequences.
All binary auxiliary proofs used by this article are supplied in the
separate \emph{Binary Foundations and Auxiliary Proofs} supplement.
Cross-document theorem references carry the prefix S and are resolved
from its current source, not from a historical PDF.  The primary article
and supplement together are a self-contained proof package; the
independent structural paper is not a premise of any measurement theorem.

\subsection{Logical prerequisites of the primary article}
The covariance factorization, variance identity and full-body path
upper are proved in Section~\ref{sec:covariance83}; the finite
Bernoulli lower is proved in Section~\ref{sec:product85}.  Together
with the interface above, these suffice for the support and
finite-discrimination results of
Sections~\ref{sec:support84}--\ref{sec:controlstability85}.
The independently complete structural paper is never a premise.
The finite-neighborhood and one-sided results of
Sections~\ref{sec:tangentneighborhood86}--\ref{sec:independent86}
use these same premises; the product converse adds only standard
fidelity properties stated in its proof.  Sections~\ref{sec:nonempty87}
and~\ref{sec:width87} give explicit nonemptiness and the block-width
law using the same primary-text logical-channel calculation.
Section~\ref{sec:reset88} adds classical feedback, receiver memory
and adaptive reset blocks through a conditional fidelity argument.
\begin{center}
\begin{tabular}{@{}>{\raggedright\arraybackslash}p{.31\textwidth}>{\raggedright\arraybackslash}p{.61\textwidth}@{}}
\toprule
Primary result & Required auxiliary result\\
\midrule
Support kernel, orbit, curve and finite-neighborhood classification &
None from the supplement: all probability, covariance and correction
premises are proved in the primary.\\
Independent product probes &
Primary normalized-factor remainder and the explicitly stated
fidelity purification, tensorization and trace inequalities.\\
Nonempty tubes and entangled blocks &
Primary Schur-complement realization, normalized-factor remainder,
corrected one-call channel, and finite Bernoulli lemma.\\
Transported regularization &
Primary covariance Jensen identity and variational duality.\\
Balanced-interior common learning &
Supplement Theorem~\ref{thm:rationallearner81} and
Corollary~\ref{cor:operatorconfidence80}; the former imports the
credited binary tomography primitive.\\
Noisy-projector comparison &
Supplement Theorem~\ref{thm:matrixmetric76} for the specified binary
pair. It is not a premise of the smooth-curve theorem.\\
Other binary boundary, coding and control proofs &
Preserved in the current supplement; not premises of the new
support or smooth-curve classification.\\
\bottomrule
\end{tabular}
\end{center}
